\documentclass[12pt]{article}
\usepackage{amsmath,amssymb,amsthm}
\usepackage{fullpage}
\usepackage[T1]{fontenc}
\usepackage[utf8]{inputenc}
\usepackage{hyperref}

\newtheorem{theorem}{Theorem}
\newtheorem{lemma}[theorem]{Lemma}
\newtheorem{corollary}[theorem]{Corollary}
\newtheorem{proposition}[theorem]{Proposition}
\theoremstyle{definition}
\newtheorem{definition}[theorem]{Definition}
\theoremstyle{remark}
\newtheorem{remark}[theorem]{Remark}

\DeclareMathOperator{\vecop}{vec}
\DeclareMathOperator{\diag}{diag}

\title{Solution to Problem 10:\\
Preconditioned Conjugate Gradient for RKHS-Constrained\\
CP Tensor Decomposition}
\author{}
\date{April 2026}

\begin{document}
\maketitle

\begin{abstract}
We describe how preconditioned conjugate gradient (PCG) solves the mode-$k$
subproblem in RKHS-constrained CP decomposition efficiently.  The system
matrix is $(Z\otimes K)^T S S^T (Z\otimes K) + \lambda(I_r\otimes K)$, of size
$nr\times nr$.  Direct solve costs $O(n^3r^3)$; we reduce this to
$O(q r^2 n + n^3 r)$ per CG iteration using structured matrix-vector products
that never form the full matrix and avoid any $O(N)$ computation.  As
preconditioner we use $P = \lambda(I_r\otimes K) + \mu I_{nr}$
(Tikhonov-type), which is diagonalisable via the eigendecomposition of $K$
at cost $O(n^3)$ once, with each preconditioner application costing $O(nr)$.
\end{abstract}

%------------------------------------------------------------------
\section{Problem Formulation}

We restate the linear system.  The unknown is $\vecop(W)\in\mathbb{R}^{nr}$;
denote $w = \vecop(W)$.  The system is
\begin{equation}\label{eq:system}
  \mathcal{A} w = b,
\end{equation}
where
\begin{align}
  \mathcal{A} &:= (Z\otimes K)^T S S^T (Z\otimes K) + \lambda(I_r\otimes K),
  \label{eq:A}\\
  b &:= (I_r\otimes K)\vecop(B). \label{eq:b}
\end{align}

Here:
\begin{itemize}
\item $K\in\mathbb{R}^{n\times n}$: symmetric PSD RKHS kernel, $n\ll M$.
\item $Z = A_d\odot\cdots\odot A_1\in\mathbb{R}^{M\times r}$: Khatri-Rao product.
\item $S\in\mathbb{R}^{N\times q}$: selection matrix ($q$ observed entries).
\item $B = TZ\in\mathbb{R}^{n\times r}$: MTTKRP (mode-$k$ tensor times
  Khatri-Rao product).
\item $\lambda > 0$: regularisation parameter.
\end{itemize}

\textbf{Key constraints}: $n,r < q \ll N$. We must not compute anything of
size $N$.

%------------------------------------------------------------------
\section{Conjugate Gradient Method}

\subsection{Applicability}

\begin{lemma}
The matrix $\mathcal{A}$ defined in~\eqref{eq:A} is symmetric positive definite
(SPD) for $\lambda > 0$ and $K\succ 0$.
\end{lemma}

\begin{proof}
$(Z\otimes K)^T S S^T (Z\otimes K) \succeq 0$ (PSD, as it is $C^T C$ for
$C = S^T(Z\otimes K)$).  The term $\lambda(I_r\otimes K)$ is SPD since
$K\succ 0$ and $\lambda > 0$.  Their sum is SPD.
\end{proof}

Since $\mathcal{A}$ is SPD, the conjugate gradient (CG) method applies and
converges in at most $nr$ iterations.  In practice, with a good preconditioner,
convergence occurs in $O(1)$ or $O(\log(1/\epsilon))$ iterations for desired
accuracy $\epsilon$.

\subsection{CG iteration cost}

Each CG iteration requires one matrix-vector product $v\mapsto\mathcal{A}v$
and one preconditioner solve $v\mapsto P^{-1}v$.  We describe both efficiently.

%------------------------------------------------------------------
\section{Efficient Matrix-Vector Product}

Given a vector $w\in\mathbb{R}^{nr}$, we compute $\mathcal{A}w$ without forming
$\mathcal{A}$ explicitly and without any $O(N)$ computation.

\subsection{Reshape}

Write $w = \vecop(W)$ with $W\in\mathbb{R}^{n\times r}$.

\subsection{First term: $(Z\otimes K)^T S S^T (Z\otimes K) w$}

\textbf{Substep 1}: Compute $v_1 = (Z\otimes K)w = (Z\otimes K)\vecop(W)$.

By the Kronecker-product identity $(A\otimes B)\vecop(X) = \vecop(BXA^T)$:
\[
  (Z\otimes K)\vecop(W) = \vecop(KWZ^T).
\]
This requires: (a) $WZ^T\in\mathbb{R}^{n\times M}$ — cost $O(nrM)$, but
$M$ is large.  \emph{We avoid this} by using $S$ as follows.

Instead, note that $S^T\vecop(V)$ (for any $V\in\mathbb{R}^{n\times M}$)
selects $q$ entries of $V$, and we only need $S^T v_1 = S^T\vecop(KWZ^T)$.

The entry at observed position $(\ell)$ of the mode-$k$ unfolding indexed
by $(i,m)$ is $(KWZ^T)_{i,m}$.  Let $\Omega\subset[n]\times[M]$ denote
the index set of observed entries (so $|\Omega|=q$).  Then:
\[
  [S^T\vecop(KWZ^T)]_\ell = (KWZ^T)_{i_\ell,m_\ell}, \quad \ell=1,\dots,q,
\]
where $(i_\ell,m_\ell)$ is the $\ell$-th observed index.

We compute this as follows:
\begin{enumerate}
\item Compute $KW\in\mathbb{R}^{n\times r}$: cost $O(n^2r)$.
\item For each observed entry $\ell$, compute $(KW)_{i_\ell,:}\cdot Z_{m_\ell,:}^T$
  (a dot product of two $r$-vectors): cost $O(qr)$ total.
\end{enumerate}
This gives $c := S^T\vecop(KWZ^T)\in\mathbb{R}^q$ at total cost $O(n^2r + qr)$.

\textbf{Substep 2}: Compute $v_2 = SS^Tv_1 = Sc\in\mathbb{R}^N$.
We do \emph{not} form $v_2$ as a dense $N$-vector.  Instead, we represent
it implicitly: $v_2$ is the sparse vector supported on $\Omega$ with values $c$.

\textbf{Substep 3}: Compute $(Z\otimes K)^T v_2 = (Z\otimes K)^T Sc$.

$(Z\otimes K)^T = (Z^T\otimes K^T) = (Z^T\otimes K)$ (since $K$ is symmetric).
We need $(Z^T\otimes K)\vecop(V)$ where $V\in\mathbb{R}^{n\times M}$ has
nonzero entries only at positions $\Omega$ (from $Sc$).

By the Kronecker identity:
$(Z^T\otimes K)\vecop(V) = \vecop(KVZ)$.

Since $V$ is sparse (only $q$ nonzero entries):
\begin{enumerate}
\item For each observed $(i_\ell,m_\ell)\in\Omega$, the column $m_\ell$ of
  $V$ has one nonzero: $V_{i_\ell,m_\ell} = c_\ell$.
\item $KVZ$: first compute $VZ\in\mathbb{R}^{n\times r}$ using sparsity:
  $(VZ)_{i,:} = \sum_\ell c_\ell\cdot\mathbf{1}[i_\ell=i]\cdot Z_{m_\ell,:}$.
  Cost: $O(qr)$ (one pass over observed entries).
\item Then $K(VZ)\in\mathbb{R}^{n\times r}$: cost $O(n^2r)$.
\end{enumerate}
Total cost for Substep 3: $O(qr + n^2r)$.

\textbf{First-term total cost}: $O(n^2r + qr)$ per matvec.

\subsection{Second term: $\lambda(I_r\otimes K)w$}

$(I_r\otimes K)\vecop(W) = \vecop(KWI_r^T) = \vecop(KW)$.
Cost: $O(n^2r)$.

\subsection{Right-hand side $b$}

$b = (I_r\otimes K)\vecop(B) = \vecop(KB)$.  This is computed once:
$B = TZ$ is the MTTKRP (cost $O(qr)$ using sparsity of $T$), then $KB$
costs $O(n^2r)$.

\subsection{Summary of matvec cost}

Per CG iteration: $\mathcal{A}w$ costs $\mathbf{O(n^2r + qr)}$.

%------------------------------------------------------------------
\section{Preconditioner}

\subsection{Choice}

We use the \emph{Tikhonov (spectral) preconditioner}:
\begin{equation}\label{eq:prec}
  P := \lambda(I_r\otimes K).
\end{equation}

\begin{remark}
A richer choice incorporating the dominant part of the data term:
$P = \mu(I_r\otimes K)$ with $\mu = \lambda + \|Z\|_F^2/n$ (approximating
the data term by its trace).  We use~\eqref{eq:prec} for simplicity.
\end{remark}

\subsection{Preconditioning solve: $P^{-1}v$}

Since $P = \lambda(I_r\otimes K)$, we have
$P^{-1} = \lambda^{-1}(I_r\otimes K^{-1})$
(assuming $K$ is invertible; otherwise use the pseudoinverse).

\textbf{Precomputation} (once): Compute the eigendecomposition $K = U\Lambda U^T$
(where $U\in O(n)$, $\Lambda = \diag(\lambda_1,\dots,\lambda_n)$, $\lambda_i>0$).
Cost: $O(n^3)$.

\textbf{Apply $P^{-1}$} to $v = \vecop(V)$:
\[
  P^{-1}v = \lambda^{-1}(I_r\otimes K^{-1})\vecop(V)
  = \lambda^{-1}\vecop(K^{-1}V).
\]
Compute $K^{-1}V = U\Lambda^{-1}U^T V$:
\begin{enumerate}
\item $U^TV$: cost $O(n^2r)$.
\item Scale rows by $\lambda_i^{-1}$: cost $O(nr)$.
\item Multiply by $U$: cost $O(n^2r)$.
\end{enumerate}
Total per application: $O(n^2r)$.

\subsection{Why this preconditioner works}

The system matrix is $\mathcal{A} = D + P$ where
$D = (Z\otimes K)^T SS^T(Z\otimes K) \succeq 0$ is the data term and
$P = \lambda(I_r\otimes K)$ is the regulariser.
The preconditioned system $P^{-1}\mathcal{A} = P^{-1}D + I$ has eigenvalues
in $[1, 1 + \|P^{-1}D\|_2]$.

Key: $\|P^{-1}D\|_2 \leq \lambda^{-1}\|(I_r\otimes K^{-1})\|_2
\cdot\|(Z\otimes K)^T SS^T(Z\otimes K)\|_2$.
For $\lambda$ large relative to the spectral norm of $D$, the condition
number $\kappa(P^{-1}\mathcal{A}) \approx 1 + O(\lambda^{-1})$, giving
rapid CG convergence.  For moderate $\lambda$, convergence in
$O(\kappa^{1/2}\log(1/\epsilon))$ iterations.

%------------------------------------------------------------------
\section{Full Complexity Analysis}

\begin{theorem}\label{thm:complexity}
Let $t_{\mathrm{CG}}$ be the number of CG iterations to reach accuracy
$\epsilon$.  The total cost of solving~\eqref{eq:system} via PCG is:
\[
  O\!\left(n^3 + t_{\mathrm{CG}}\cdot(n^2r + qr)\right).
\]
Under mild spectral assumptions on $K$ and the data (e.g.,
$\kappa(K)\lesssim 1$), $t_{\mathrm{CG}} = O(\log(1/\epsilon))$, giving
total cost $\mathbf{O(n^3 + (n^2 r + qr)\log(1/\epsilon))}$.
\end{theorem}

\begin{proof}
\begin{itemize}
\item Precomputation: eigendecomposition of $K$ in $O(n^3)$.
  Computation of $KB\in\mathbb{R}^{n\times r}$ in $O(n^2r + qr)$.
\item Per iteration: matvec $O(n^2r+qr)$, preconditioner solve $O(n^2r)$,
  CG updates $O(nr)$.  Total per iteration: $O(n^2r + qr)$.
\item Number of iterations: $O(\kappa^{1/2}\log(1/\epsilon))$ where
  $\kappa = \kappa(P^{-1}\mathcal{A})$.
\end{itemize}
All computations involve arrays of size at most $\max(nr,qr,n^2) \ll N$.
No $O(N)$ operation is performed.
\end{proof}

\noindent\textbf{Comparison with direct solve.}
Direct Cholesky factorisation of $\mathcal{A}$ costs $O(n^3r^3)$, which is
much larger than $O(n^3 + qr\log(1/\epsilon))$ under the assumption $n,r<q\ll N$.
For example, $n=100$, $r=10$, $q=10^4$: direct solve costs $\approx 10^9$
operations; PCG costs $\approx 10^4\cdot 10^3 + 10^6 = 10^7 + 10^6$
operations.

%------------------------------------------------------------------
\section{Summary of the Algorithm}

\begin{enumerate}
\item \textbf{Precompute}: $K = U\Lambda U^T$ ($O(n^3)$);
  $B = TZ$ ($O(qr)$); $b = \vecop(KB)$ ($O(n^2r)$).
\item \textbf{Initialise}: $w_0 = 0$, $r_0 = b$, $p_0 = P^{-1}r_0$.
\item \textbf{CG loop} (for $j = 0,1,2,\dots$):
  \begin{enumerate}
  \item $q_j = \mathcal{A}p_j$ (matvec via structured formula, $O(n^2r+qr)$)
  \item $\alpha_j = (r_j^T P^{-1}r_j)/(p_j^T q_j)$
  \item $w_{j+1} = w_j + \alpha_j p_j$
  \item $r_{j+1} = r_j - \alpha_j q_j$
  \item $z_{j+1} = P^{-1}r_{j+1}$ (preconditioner solve, $O(n^2r)$)
  \item $\beta_j = (r_{j+1}^T z_{j+1})/(r_j^T P^{-1}r_j)$
  \item $p_{j+1} = z_{j+1} + \beta_j p_j$
  \end{enumerate}
\item \textbf{Output}: $W = \mathrm{reshape}(w, n, r)$.
\end{enumerate}

%------------------------------------------------------------------
\section*{References}

\begin{itemize}
\item Y.~Saad, \emph{Iterative Methods for Sparse Linear Systems}, 2nd ed.,
  SIAM, 2003.
\item T.~G.~Kolda and B.~W.~Bader, ``Tensor decompositions and applications,''
  \emph{SIAM Review}, 51(3):455--500, 2009.
\item C.~A.~Micchelli and M.~Pontil, ``On learning vector-valued functions,''
  \emph{Neural Computation}, 17(1):177--204, 2005.
\item N.~Halko, P.~G.~Martinsson, and J.~A.~Tropp, ``Finding structure with
  randomness: Probabilistic algorithms for constructing approximate matrix
  decompositions,'' \emph{SIAM Review}, 53(2):217--288, 2011.
\end{itemize}

\end{document}
